\section{Analytic classes and changes of physical coordinates}\label{sec:model}
The geometric constructions prescribe a smooth velocity first. This section
records exactly which other solutions are excluded by the energy argument.
Pressure assumptions are kept separate: continuity in a pressure norm, a
norm at each individual time, and an integrable pressure gradient are
different conditions.

Write $\mathcal F_\nu[U]=U_t+(U\cdot\nabla)U-\nu\Delta U$.
All noncompact results below use $\mathbb R^3$; their velocities and forces
have one fixed compact spatial support unless the individual statement
says otherwise. This property concerns the constructed solution, not a
competing solution. On $[0,T]$, write $CH^k=C([0,T];H^k(\mathbb R^3))$
and $C^1L^2=C^1([0,T];L^2(\mathbb R^3))$. Spatial constants in pressure
may depend on time.

\begin{lemma}[The pressure term without a pressure norm]\label{lem:b-gradient}
If $w,g\in L^2(\mathbb R^3;\mathbb R^3)$ satisfy
$\operatorname{div}w=0$ and $\operatorname{curl}g=0$ in distributions, then
$\langle w,g\rangle_{L^2}=0$. In particular, any distributional gradient
$g=\nabla p\in L^2$ has zero pairing with $w$, without an assumption that
$p\in L^2$.
\end{lemma}
\begin{proof}
An $L^2$ field is a tempered distribution. The differential identities,
initially tested on compactly supported smooth functions, extend to
Schwartz tests by multiplying by cutoffs tending to one and using $L^2$
convergence of the tests and their first derivatives. Fourier transformation
gives $\xi\cdot\widehat w=0$ and
$\xi_j\widehat g_k-\xi_k\widehat g_j=0$. These are identities of locally
integrable functions, so they hold almost everywhere. Away from $\xi=0$,
$\widehat g$ is parallel to $\xi$ and $\widehat w$ is perpendicular to it.
Their Hermitian product is zero. The exceptional point has measure zero;
Plancherel proves the assertion.
\end{proof}

\begin{lemma}[Residual realization and separate comparison classes]
\label{lem:b-comparison}\label{lem:p2-realization}
Let $\nu>0$ and let $U$ be a prescribed smooth divergence-free velocity
on $[0,\infty)\times\mathbb R^3$, with $U(0)=0$. On each finite interval,
assume $U\in CH^2\cap C^1L^2$ and $U,\nabla U$ are bounded. Put
$f=\mathcal F_\nu[U]$. Then $(U,0)$ is a solution. Its velocity is unique
in each of the following separately stated comparison classes, with the
same force and initial data:
\begin{enumerate}[label=(\roman*),itemsep=2pt]
\item Smooth classical $v\in CH^2\cap C^1L^2$, with bounded $v$, and a
pressure representative $p\in CH^1$.
\item Smooth classical $v\in CH^2\cap C^1L^2$, with bounded $v,\nabla v$,
and an $H^1$ pressure representative at each time, with no required time
continuity in that pressure norm.
\item Strong solutions with $v\in CH^4\cap C^1H^2$ and distributional
$\nabla p\in CL^2$.
\item Smooth classical $v\in CH^2\cap C^1L^2$, with bounded $v,\nabla v$,
and a pressure representative $p\in CL^2$.
\item Smooth classical $v\in CH^2\cap C^1L^2$, with bounded $v,\nabla v$,
without an additional integrability or normalization condition on pressure.
\end{enumerate}
The pressure gradient is unique. An $L^2$ pressure representative, when
specified, is zero; otherwise pressure is determined up to a function of
time. The reference material flow is a smooth diffeomorphism at every
finite time and exists for every finite forward time.

On a flat torus $\T_L^3$, the same conclusion holds in the classical class
where $v,v_t$, spatial derivatives through order two, and $p,\nabla p$ are
continuous on finite closed cylinders, with periodic mean-zero pressure.
No Sobolev conditions at infinity are involved in that assertion.
\end{lemma}
\begin{proof}
Substitution proves existence of the prescribed solution. Set $w=v-U$.
The difference equation is
\[
 w_t+(v\cdot\nabla)w+(w\cdot\nabla)U
    =\nu\Delta w-\nabla p,\qquad \operatorname{div}w=0.
\]
The stated velocity regularity implies $w\in C^1L^2\cap CH^2$, so
$\frac d{dt}\|w\|_2^2=2\langle w,w_t\rangle$ at each time. In case
(iii), the stronger assumptions imply these inclusions. They also imply
bounded velocity and spatial derivatives through order two: for $j\le2$,
Cauchy--Schwarz in Fourier space uses the finite integral
$\int_{\mathbb R^3}|\xi|^{2j}(1+|\xi|^2)^{-4}\,d\xi$.

For the transport term, insert a cutoff $\chi_R$ that equals one on the
ball of radius $R$, is supported in the ball of radius $2R$, and satisfies
$|\nabla\chi_R|\le C/R$. Its boundary error is at most
$CR^{-1}\|v\|_\infty\|w\|_2^2$, which tends to zero. Diffusion is
integrated by parts in $H^2$; alternatively its cutoff error is bounded by
$CR^{-1}\|w\|_2\|\nabla w\|_2$. The deformation term is bounded by
$\|\nabla U\|_{\infty,\mathrm{op}}\|w\|_2^2$.

The same pressure cancellation applies in all five classes. At each time,
the equation gives the following equality of distributional gradients in
$L^2$:
\[
 \nabla p=f-v_t-(v\cdot\nabla)v+\nu\Delta v\in L^2,
 \qquad \|(v\cdot\nabla)v\|_2\le\|v\|_\infty\|\nabla v\|_2.
\]
The velocity assumptions in every class give all terms on the right in
$L^2$; the identical estimate for $U$ and its stated time and space
regularity give $f\in L^2$. Since a distributional gradient is curl free,
Lemma~\ref{lem:b-gradient} gives $\langle\nabla p,w\rangle=0$.
This cancellation neither chooses an $L^2$ pressure representative nor
uses continuity of pressure in a norm.

In every case the resulting pointwise energy inequality is
\[
 \frac12\frac d{dt}\|w\|_2^2+\nu\|\nabla w\|_2^2
 \le \|\nabla U\|_{\infty,\mathrm{op}}\|w\|_2^2.
\]
The coefficient is bounded on each finite interval and $w(0)=0$, so an
integrating factor gives $w=0$. The pressure has already disappeared before
time integration; the proof imposes no unlisted temporal pressure norm.
Substitution then gives $\nabla p=0$. A spatially constant $L^2$ function
on $\mathbb R^3$ is zero, which proves the normalization assertion.

On the torus every integration by parts is periodic and has no boundary
term. The specified classical regularity justifies norm differentiation,
diffusion, and pressure cancellation. The same energy inequality gives
uniqueness and the mean-zero normalization fixes pressure. In either domain,
the smooth reference velocity and its bounded spatial derivative give
unique material trajectories; bounded speed prevents finite-time escape.
Solving backwards on a finite interval gives the inverse flow, and smooth
ODE dependence gives its smoothness.
\end{proof}

Cases (i)--(v) identify the precise hypotheses used by the applications;
their presence in a common lemma does not replace a theorem's declared
comparison class by a different one. In particular, a pressure norm is not
silently imposed in a gradient-only application. The proof is the classical
energy comparison for prescribed smooth solutions \cite[Section~18]{Leray1934},
with its noncompact pressure pairings made explicit.

For the flow $\Phi_t$, differentiation of its trajectory equation gives
\[
 U(t)=(\partial_t\Phi_t)\circ\Phi_t^{-1},\qquad
 \partial_{tt}\Phi_t(a)=
 (f-\nabla p+\nu\Delta U)(t,\Phi_t(a)).
\]
These identities describe the material form of the same constructed solution;
they do not add an independent existence claim.

\begin{lemma}[Physical scaling]\label{lem:p2-chart}
For $a,b>0$, $x=x_0+ay$, and $s=bt$, let
$U(t,x)=abV(bt,(x-x_0)/a)$. Then trajectories are transformed by this change
of variables, incompressibility is preserved, and at the prescribed physical
viscosity $\nu$,
\[
 \mathcal F_\nu[U](t,x)=ab^2\left[
 V_s+(V\cdot\nabla_y)V-\frac{\nu}{a^2b}\Delta_yV
 \right](bt,(x-x_0)/a).
\]
\end{lemma}
\begin{proof}
The trajectory equation gives $\dot x=abV$; differentiating in space gives
$\operatorname{div}_xU=b\operatorname{div}_yV$. The three residual terms
scale by $ab^2,ab^2$, and $b/a$, respectively, proving the formula.
\end{proof}
Thus a side-$10$ torus is kept as a physical side-$10$ torus. If a chart
is rescaled, the residual is recomputed with the fixed $\nu$; the equation
is not identified with a differently normalized viscosity.

\begin{proposition}[Projection on a flat torus]\label{prop:projection-l2}
An effective smooth mean-zero force $g$ on $\T_L^3$ has an effective
mean-zero potential $\phi$ satisfying $\Delta\phi=\operatorname{div}g$.
Replacing $(g,p)$ by $(g-\nabla\phi,p-\phi)$ preserves velocity and makes
the force solenoidal. Uniform mixed-derivative bounds, temporal $L^2$
bounds of their spatial suprema, and separately imposed time periodicity,
eventual stationarity, or orderwise decay bounds are preserved.
\end{proposition}
\begin{proof}
In frequencies $\xi_k=2\pi k/L$, set
$\widehat\phi(k)=-i\xi_k\cdot\widehat g(k)/|\xi_k|^2$ for $k\ne0$
and zero for $k=0$. The multiplier from $g$ to $\nabla\phi$ is
$\xi_k\xi_k^{\mathsf T}/|\xi_k|^2$. For a spatial derivative order $r$,
integrate each Fourier coefficient $r+5$ times in a coordinate with largest
$|k_i|$. There are $O(n^2)$ lattice points on the shell $\|k\|_\infty=n$,
so the derivative series and its tail converge absolutely, with
\[
 \|\partial_t^h\partial_x^\alpha\nabla\phi\|_\infty
 \le C_{\alpha,L}\max_i
       \|\partial_t^h\partial_{x_i}^{|\alpha|+5}g\|_\infty.
\]
This bound is pointwise in time, and proves each asserted norm or weighted
time estimate. Derivative bounds give effective Fourier tails; coefficient
integrals can be computed by Riemann sums with derivative error estimates.
The formula gives the Poisson identity, divergence cancellation, and the
pressure sign by substitution. Linearity preserves the time symmetries.
\end{proof}
Projection generally changes pressure. It supplies no compact-support
conclusion on Euclidean space and is distinct from the direct zero-pressure
solenoidal shear construction.

\paragraph{Effective flat profiles.}
For later cutoff formulas put
\begin{equation}\label{eq:p2-step}
 \rho(s)=\begin{cases}e^{-1/s},&s>0,\\0,&s\le0,\end{cases}
 \qquad \sigma(s)=\frac{\rho(s)}{\rho(s)+\rho(1-s)}.
\end{equation}
Positive-side derivatives of $\rho$ are polynomials in $1/s$ times
$e^{-1/s}$ and extend flatly at zero. The estimate
$v^me^{-v}\le(m+k)!v^{-k}$ gives effective error bounds near the seam,
and $\rho(s)+\rho(1-s)\ge e^{-2}$. Products, translations, rescalings,
and derivatives therefore yield effective cutoffs with all requested
derivative bounds. Periodization uses zero collars; at approximate real
arguments a finite superset of possible translates is evaluated. No real
equality test at a seam or execution of the encoded computation is used.
